\section{Chi Tiết Triển Khai}

\subsection{Jenkins Infrastructure \& Pipeline}

\subsubsection{Cài Đặt Jenkins Server}

\paragraph{Môi Trường Triển Khai}

\begin{table}[H]
\centering
\caption{Thông số môi trường Jenkins}
\begin{tabular}{|l|p{9cm}|}
\hline
\textbf{Thông số} & \textbf{Giá trị} \\
\hline
Phương thức triển khai & Virtual Machine Amazon EC2 \\
\hline
Phiên bản Jenkins & 2.555.1 \\
\hline
Hệ điều hành & Amazon Linux 2023 (kernel 6.1 AMI, architecture: arm64, instance: t3.small) \\
\hline
Phiên bản Java & Amazon Corretto 21 (OpenJDK 21.0.10 LTS) \\
\hline
\end{tabular}
\end{table}

\paragraph{Plugin Đã Cài Đặt}

\begin{table}[H]
\centering
\caption{Các plugin Jenkins đã cài đặt}
\begin{tabular}{|l|l|}
\hline
\textbf{Plugin} & \textbf{Mục đích} \\
\hline
Git Plugin & Kết nối với GitHub repository \\
\hline
Pipeline & Chạy Jenkinsfile dạng Declarative/Scripted \\
\hline
Multibranch Pipeline & Tự động quét và build nhiều branch \\
\hline
JaCoCo Plugin & Publish báo cáo độ phủ unit test \\
\hline
Warnings Next Generation & Hiển thị kết quả phân tích tĩnh \\
\hline
GitHub Integration Plugin & Nhận sự kiện từ GitHub Webhook \\
\hline
\end{tabular}
\end{table}

Ngoài ra còn một số plugin được Jenkins recommend khi setup lần đầu tiên.

\begin{figure}[H]
\centering
\reportimg{jenkins/jenkins-dashboard.jpg}
\caption{Jenkins Dashboard sau khi cài đặt thành công}
\end{figure}

\subsubsection{Cấu Hình GitHub Webhook}

Webhook được tạo tại: \texttt{Repository > Settings > Webhooks > Add webhook}

\begin{table}[H]
\centering
\caption{Cấu hình GitHub Webhook}
\begin{tabular}{|l|l|}
\hline
\textbf{Cấu hình} & \textbf{Giá trị} \\
\hline
Payload URL & \texttt{http://18.143.92.157:8080/github-webhook/} \\
\hline
Content type & \texttt{application/json} \\
\hline
Trigger events & Push events, Pull request events \\
\hline
Trạng thái & Active \\
\hline
\end{tabular}
\end{table}

\begin{figure}[H]
\centering
\reportimg{jenkins/webhook-config.jpg}
\caption{Cấu hình Webhook trên GitHub}
\end{figure}

\begin{figure}[H]
\centering
\reportimg{jenkins/webhook-pr-trigger.jpg}
\caption{Cấu hình trigger event Pull Request cho Webhook}
\end{figure}

\begin{figure}[H]
\centering
\reportimg{jenkins/webhook-ping-success.jpg}
\caption{Webhook delivery thành công (ping event trả về HTTP 200)}
\end{figure}

\subsubsection{Thiết Lập Multibranch Pipeline}

\begin{table}[H]
\centering
\caption{Cấu hình Multibranch Pipeline Job}
\begin{tabular}{|l|l|}
\hline
\textbf{Cấu hình} & \textbf{Giá trị} \\
\hline
Tên job & YAS \\
\hline
Loại job & Multibranch Pipeline \\
\hline
Branch Source & GitHub \\
\hline
URL Repository & \url{https://github.com/com-suon-bi-cha/yas.git} \\
\hline
Scan interval & 2 minute \\
\hline
Đường dẫn Jenkinsfile & \texttt{Jenkinsfile} (tại root) \\
\hline
\end{tabular}
\end{table}

\begin{figure}[H]
\centering
\reportimg{jenkins/multibranch-config-1.jpg}
\caption{Cấu hình Multibranch Pipeline Job (1)}
\end{figure}

\begin{figure}[H]
\centering
\reportimg{jenkins/multibranch-config-2.jpg}
\caption{Cấu hình Multibranch Pipeline Job (2)}
\end{figure}

\begin{figure}[H]
\centering
\reportimg{jenkins/multibranch-config-pat.jpg}
\caption{Cấu hình Personal Access Token cho Multibranch Pipeline}
\end{figure}

\begin{figure}[H]
\centering
\reportimg{jenkins/multibranch-scan-branches.jpg}
\caption{Multibranch Pipeline — danh sách branch được Jenkins phát hiện}
\end{figure}

\begin{figure}[H]
\centering
\reportimg{jenkins/pipeline-auto-trigger.jpg}
\caption{Pipeline tự động kích hoạt sau khi push code}
\end{figure}

\subsubsection{Nội Dung Jenkinsfile}

\paragraph{Cấu Trúc Pipeline}

Pipeline gồm 7 stages theo thứ tự:

\begin{lstlisting}[language=Java, caption={Cấu trúc Jenkinsfile}]
pipeline {
    agent { label 'individual-agent' }

    stages {
        stage('Pre-check')          { ... }
        stage('Secret Scanning')    { ... }
        stage('Monorepo Execution') { ... }
        stage('Code Quality')       { ... }
        stage('Quality Gate')       { ... }
        stage('Coverage Report')    { ... }
        stage('Dependency Scan')    { ... }
    }
}
\end{lstlisting}

\paragraph{Logic Monorepo Execution (Test $\rightarrow$ Upload $\rightarrow$ Build)}

Đây là stage cốt lõi, thực hiện cả 2 giai đoạn \textbf{Test} và \textbf{Build}. Pipeline sử dụng \texttt{git diff} để chỉ xử lý service có thay đổi:

\begin{lstlisting}[language=Java, caption={Logic Monorepo Execution trong Jenkinsfile}]
stage('Monorepo Execution') {
    steps {
        script {
            def changedFiles = sh(
                script: 'git diff --name-only HEAD~1 HEAD',
                returnStdout: true
            ).trim().split('\n')

            def services = [
                'media', 'product', 'order', 'inventory',
                'payment', 'promotion', 'rating', 'delivery',
                'sampledata', 'recommendation', 'customer',
                'location', 'cart', 'tax', 'search', 'webhook',
                'common-library', 'backoffice-bff', 'storefront-bff'
            ]

            for (service in services) {
                if (changedFiles.any { it.startsWith("${service}/") }) {
                    sh "mvn test -pl ${service} -am"             // Test
                    junit testResults: "${service}/target/surefire-reports/*.xml",
                          allowEmptyResults: true                // Upload
                    sh "mvn package -DskipTests -pl ${service} -am" // Build
                }
            }
        }
    }
}
\end{lstlisting}

\textbf{Giải thích 3 bước:}
\begin{itemize}
    \item \texttt{mvn test} — chạy toàn bộ unit test của service, JaCoCo agent tự động thu thập coverage data.
    \item \texttt{junit} — upload kết quả test (file XML) lên Jenkins UI để hiển thị trong Test Result.
    \item \texttt{mvn package -DskipTests} — build artifact \texttt{.jar} mà không chạy lại test, tiết kiệm thời gian.
\end{itemize}

\subsubsection{Kiểm Tra Pipeline Hoạt Động}

\begin{figure}[H]
\centering
\reportimg{jenkins/pipeline-all-stages-success.jpg}
\caption{Tất cả stage pipeline chạy thành công}
\end{figure}

\begin{figure}[H]
\centering
\reportimg{jenkins/pipeline-build-history.jpg}
\caption{Lịch sử build trên Jenkins}
\end{figure}

\subsubsection{Cấu Hình Agent Pool (Distributed Builds)}

Master node (Built-In Node) của Jenkins chỉ có \textbf{2 GB RAM} (instance \texttt{t3.small}), không đủ tài nguyên để build/test các service Java lớn trong monorepo YAS. Nhóm triển khai mô hình \textbf{Distributed Builds} với một pool gồm 4 agent riêng biệt.

\textbf{Lợi ích:}
\begin{itemize}
    \item \textbf{Giảm tải master node:} Master chỉ đảm nhiệm điều phối (scheduling), toàn bộ workload build/test được chuyển sang agent.
    \item \textbf{Tận dụng tài nguyên cá nhân:} Mỗi thành viên cấu hình máy cá nhân làm agent, chủ động bật agent khi có thay đổi cần chạy pipeline.
    \item \textbf{Song song hóa:} Nhiều pipeline có thể chạy đồng thời trên các agent khác nhau.
    \item \textbf{Môi trường linh hoạt:} Agent có thể chạy trên nhiều kiến trúc khác nhau (aarch64, amd64).
\end{itemize}

\textbf{Danh sách agent:}
\begin{itemize}
    \item \texttt{anamap-agent}
    \item \texttt{bingsu-agent}
    \item \texttt{nhatdat-agent}
    \item \texttt{tunah-agent}
\end{itemize}

Tất cả agent được gán chung label \texttt{individual-agent} để Jenkins có thể tự động phân phối job đến bất kỳ agent nào đang online.

\begin{lstlisting}[language=Java, caption={Cấu hình agent label trong Jenkinsfile}]
pipeline {
    agent { label 'individual-agent' }
    // ...
}
\end{lstlisting}

\begin{figure}[H]
\centering
\reportimg{jenkins/agent-pool-nodes.png}
\caption{Danh sách Nodes trên Jenkins: 4 agent + Built-In Node}
\end{figure}


